\documentclass[11pt]{article}
\usepackage[a4paper,margin=2.4cm]{geometry}
\usepackage{amsmath,amssymb}
\usepackage{enumitem}
\usepackage{hyperref}

\title{Lecture 4 Plan\\Basic Dependent Type Theory and an Agda Demo}
\author{}
\date{}

\begin{document}
\maketitle

\section*{Learning objectives}
By the end of this three-hour unit, students should be able to:
\begin{enumerate}[nosep]
  \item explain how dependent types extend STLC and System F;
  \item read a dependent function type $(x : A) \to B\;x$;
  \item understand indexed families through length-indexed vectors;
  \item see propositions-as-types through the equality type and \texttt{refl};
  \item use Agda's type checker as a small proof checker.
\end{enumerate}

\section*{Connection to previous units}
\begin{center}
\begin{tabular}{l l}
STLC & terms depend on terms; types classify terms.\\
System F & terms can abstract over types.\\
Dependent type theory & types may depend on terms.\\
Agda & type checking can verify programs and proofs.\\
\end{tabular}
\end{center}

\section*{Suggested pacing}
\begin{center}
\begin{tabular}{r p{0.68\linewidth}}
0--20 min & From function types to dependent function types.\\
20--45 min & Indexed families: vectors indexed by length.\\
45--65 min & Exercise 1: read vector types.\\
65--75 min & Break.\\
75--105 min & Equality as a type and proof by computation.\\
105--145 min & Live Agda demo: \texttt{Nat}, \texttt{Vec}, \texttt{append}, \texttt{head}.\\
145--170 min & Exercise 2: fill small definitions and predict rejected cases.\\
170--180 min & Summary and further directions.\\
\end{tabular}
\end{center}

\section*{Core ideas}
\subsection*{Dependent functions}
Ordinary function type:
\[
  A \to B
\]
Dependent function type:
\[
  (x : A) \to B\;x
\]
The return type may mention the input value.

\subsection*{Indexed families}
A vector type can mention its length:
\[
  \mathsf{Vec}\;A\;n
\]
This should be read as: vectors of elements of type $A$ whose length is $n$.

\subsection*{Propositions as types}
A proposition is represented by a type, and a proof is represented by an inhabitant. Equality has the canonical proof:
\[
  \mathsf{refl} : x \equiv x.
\]

\section*{Live demo}
Use \texttt{agda/DependentDemo.agda}. The demo is intentionally self-contained and does not depend on the Agda standard library.

\begin{enumerate}[nosep]
  \item Define natural numbers.
  \item Define addition.
  \item Define vectors indexed by length.
  \item Define append with type \texttt{Vec A m -> Vec A n -> Vec A (m + n)}.
  \item Define head only for non-empty vectors.
  \item Show that some equalities are proved by \texttt{refl} because both sides compute to the same normal form.
\end{enumerate}

\section*{Exercise 1: reading dependent types}
Explain what each type promises.
\begin{enumerate}
  \item \texttt{Vec Nat zero}
  \item \texttt{Vec Nat (suc zero)}
  \item \texttt{\{A : Set\} \{n : Nat\} -> Vec A (suc n) -> A}
  \item \texttt{Vec A m -> Vec A n -> Vec A (m + n)}
\end{enumerate}

\paragraph{Solution checkpoints.}
\begin{enumerate}[nosep]
  \item Empty vectors of naturals.
  \item Singleton vectors of naturals.
  \item A function that extracts an element from a statically non-empty vector.
  \item A function that appends two vectors and statically preserves the length equation.
\end{enumerate}

\section*{Exercise 2: Agda discussion prompts}
After the live demo, ask students:
\begin{enumerate}
  \item Why is there no \texttt{head []} case?
  \item Why does \texttt{append [] ys = ys} typecheck?
  \item Why does the recursive case of \texttt{append} produce \texttt{Vec A (suc (m + n))}?
  \item What does \texttt{refl} prove in \texttt{ex2}? What computation is Agda using?
\end{enumerate}

\section*{Instructor notes}
Keep the unit focused on dependent types as a practical extension of earlier typing judgments. Avoid universes, identity elimination, Cubical Agda, HoTT, and normalization proofs unless there is substantial extra time.

\end{document}
